%==============================================================================
%  k-Nearest Neighbours (KNN):
%  Background, Theory, Types and Worked Examples
%
%  Self-contained lecture notes for an introductory Machine Learning course
%  (B.Tech / M.Tech level).  Prerequisites: basic coordinate geometry.
%
%  Compile with:  pdflatex knn.tex   (run twice for refs)
%  Requires a full TeX distribution (TeX Live / MiKTeX) with pgfplots + tcolorbox.
%==============================================================================
\documentclass[11pt,a4paper]{report}

%------------------------------- Page geometry --------------------------------
\usepackage[a4paper,margin=1in]{geometry}

%------------------------------- Mathematics ----------------------------------
\usepackage{amsmath,amssymb,amsthm,mathtools}
\usepackage{bm}

%------------------------------- Tables & layout ------------------------------
\usepackage{booktabs}
\usepackage{array}
\usepackage{enumitem}
\usepackage{caption}
\usepackage{multirow}
%------------------------------- Graphics & boxes -----------------------------
\usepackage{graphicx}
\usepackage{xcolor}
\usepackage{tcolorbox}
\usepackage{tikz}
\usepackage{pgfplots}
\pgfplotsset{compat=1.18}
\usetikzlibrary{arrows.meta,calc}

%------------------------------- Code listings --------------------------------
\usepackage{listings}

%------------------------------- Hyperlinks (load last) -----------------------
\usepackage[colorlinks=true,linkcolor=blue!55!black,citecolor=blue!55!black,
urlcolor=blue!55!black,linktocpage=true]{hyperref}

%------------------------------- Palette --------------------------------------
\definecolor{cnavy}{HTML}{264653}
\definecolor{cteal}{HTML}{2A9D8F}
\definecolor{corange}{HTML}{F4A261}
\definecolor{ccoral}{HTML}{E76F51}
\definecolor{lightbg}{HTML}{F4F7F8}

%------------------------------- Boxes ----------------------------------------
\newtcolorbox{conceptbox}{
	colback=lightbg,colframe=cnavy,boxrule=0.8pt,arc=2mm,
	left=2mm,right=2mm,top=1.5mm,bottom=1.5mm}
\newtcolorbox{resultbox}{
	colback=cteal!8,colframe=cteal,boxrule=1pt,arc=2mm,
	left=2mm,right=2mm,top=1.5mm,bottom=1.5mm}
\newcommand{\classpass}{\textbf{\textcolor{cteal}{Pass}}}
\newcommand{\classfail}{\textbf{\textcolor{ccoral}{Fail}}}

%------------------------------- Theorem-like envs ----------------------------
\theoremstyle{definition}
\newtheorem{definition}{Definition}[chapter]
\newtheorem{example}{Example}[chapter]
\newtheorem{exercise}{Exercise}[chapter]
\theoremstyle{remark}
\newtheorem*{remark}{Remark}

%------------------------------- Handy macros ---------------------------------
\newcommand{\R}{\mathbb{R}}
\newcommand{\vx}{\mathbf{x}}
\newcommand{\vy}{\mathbf{y}}
\DeclareMathOperator*{\argmax}{arg\,max}
\DeclareMathOperator*{\argmin}{arg\,min}
\DeclareMathOperator{\mode}{mode}

%------------------------------- Listing style --------------------------------
\definecolor{codegreen}{rgb}{0.0,0.5,0.0}
\definecolor{codegray}{rgb}{0.5,0.5,0.5}
\definecolor{codeblue}{rgb}{0.1,0.1,0.6}
\definecolor{backcol}{rgb}{0.97,0.97,0.95}
\lstdefinestyle{python}{
	backgroundcolor=\color{backcol},language=Python,
	commentstyle=\color{codegreen},keywordstyle=\color{codeblue}\bfseries,
	numberstyle=\tiny\color{codegray},stringstyle=\color{codegreen},
	basicstyle=\ttfamily\footnotesize,breaklines=true,captionpos=b,
	numbers=left,numbersep=6pt,showstringspaces=false,frame=single,
	rulecolor=\color{codegray},tabsize=2}

%==============================================================================
\begin{document}
	
	%------------------------------- Title ----------------------------------------
	\begin{titlepage}
		\centering
		\vspace*{2cm}
		{\Huge\bfseries $k$-Nearest Neighbours (KNN)\par}
		\vspace{0.6cm}
		{\Large Background, Theory, Types\\[2pt]
			and Worked Examples by Hand\par}
		\vspace{1.5cm}
		{\large A Self-Contained Set of Lecture Notes\par}
		\vspace{0.4cm}
		{\large Introduction to Machine Learning\par}
		\vspace{2.5cm}
		\begin{tabular}{c}
			\hline\\[-6pt]
			\textbf{Prerequisites:} Coordinate Geometry and Basic Statistics\\[4pt]
			\textbf{Level:} B.Tech / M.Tech\\[4pt]
			\hline
		\end{tabular}
		\vfill
		{\large \today\par}
	\end{titlepage}
	
	\tableofcontents
	\clearpage
	
	%==============================================================================
	\chapter{Background: KNN in Machine Learning}
	%==============================================================================
	
	\section{Supervised vs.\ Unsupervised Learning}
	In \emph{supervised} learning each training example carries both input features
	and a known target, and we learn a mapping from inputs to targets. In
	\emph{unsupervised} learning there are no targets. \textbf{$k$-nearest
		neighbours (KNN) is a supervised algorithm}: it needs labelled training data.
	
	\begin{conceptbox}
		The task is supervised because the training data contains both the input
		features $\vx=(x_1,\dots,x_n)$ and the known target $y$:
		\[
		\vx=(x_1,\dots,x_n)\ \longrightarrow\ y .
		\]
		For \emph{classification} $y$ is a discrete class label (e.g.\ Pass/Fail); for
		\emph{regression} $y$ is a real number.
	\end{conceptbox}
	
	\section{Classification vs.\ Regression}
	\begin{center}
		\begin{tabular}{@{}p{6.4cm}p{6.4cm}@{}}
			\toprule
			\textbf{Classification} & \textbf{Regression}\\
			\midrule
			Target is a discrete label. & Target is a continuous number.\\
			KNN rule: \emph{majority vote} of neighbours. & KNN rule: \emph{average} of
			neighbours' values.\\
			Example: spam / not-spam. & Example: predict a house price.\\
			\bottomrule
		\end{tabular}
	\end{center}
	KNN handles \textbf{both}, changing only the final aggregation step.
	
	\section{Where does KNN Fit? Lazy, Instance-Based Learning}
	Most algorithms (linear/logistic regression, decision trees, neural networks)
	\emph{learn a model} at training time and then discard the data---they are
	\emph{eager} learners. KNN does the opposite:
	
	\begin{conceptbox}
		Basic KNN builds \textbf{no explicit model} during training. It simply
		\emph{stores} the labelled examples, and does all the work---distance
		computation, neighbour search, voting---only when a new point must be
		classified. It is therefore called a \textbf{lazy} (or \textbf{instance-based},
		or \textbf{memory-based}) learner.
		\[
		\text{new point}\rightarrow\text{distances}\rightarrow\text{$k$ nearest}
		\rightarrow\text{vote / average}\rightarrow\text{prediction.}
		\]
	\end{conceptbox}
	
	\section{Applications}
	\begin{itemize}[nosep]
		\item Recommendation systems (users/items similar to yours).
		\item Handwritten-digit and image recognition.
		\item Medical diagnosis from patient measurements.
		\item Credit scoring and anomaly detection.
		\item Imputing missing values (nearest-neighbour imputation).
	\end{itemize}
	
	%==============================================================================
	\chapter{Distance and Similarity Measures}
	%==============================================================================
	
	KNN is built entirely on a notion of ``closeness''. Let
	$\vx=(x_1,\dots,x_n)$ and $\vy=(y_1,\dots,y_n)$ be two points (two different
	observations); $x_k$ is the $k$-th feature of $\vx$.
	
	\section{Common Distance Measures}
	\begin{definition}[Euclidean distance ($L_2$)]
		\begin{equation}
			d(\vx,\vy)=\sqrt{\sum_{k=1}^{n}(x_k-y_k)^2}.
			\label{eq:euclid}
		\end{equation}
		The straight-line distance; the default for numerical features.
	\end{definition}
	\begin{definition}[Manhattan distance ($L_1$)]
		\[ d(\vx,\vy)=\sum_{k=1}^{n}\lvert x_k-y_k\rvert . \]
	\end{definition}
	\begin{definition}[Minkowski distance ($L_p$)]
		\[ d(\vx,\vy)=\Bigl(\sum_{k=1}^{n}\lvert x_k-y_k\rvert^{p}\Bigr)^{1/p}, \]
		which gives Manhattan at $p=1$ and Euclidean at $p=2$.
	\end{definition}
	\begin{definition}[Hamming distance]
		For categorical features, the number of positions at which $\vx$ and $\vy$
		differ. Used for binary/string data.
	\end{definition}
	\begin{definition}[Cosine dissimilarity]
		$d(\vx,\vy)=1-\dfrac{\vx^{\!\top}\vy}{\lVert\vx\rVert\,\lVert\vy\rVert}$; measures
		angle rather than magnitude (common for text).
	\end{definition}
	
	\section{Why Feature Scaling Matters}
	\label{sec:scaling}
	Euclidean distance is dominated by features with a large numerical range.
	
	\begin{conceptbox}
		If ``hours studied'' ranges $1$--$6$ but ``attendance'' ranges $50$--$80$, then
		in $\sqrt{(\Delta\text{hours})^2+(\Delta\text{attendance})^2}$ the attendance
		term dwarfs the hours term---KNN would effectively ignore hours. To give every
		feature a fair say we \textbf{scale} first.
	\end{conceptbox}
	
	\begin{itemize}[nosep]
		\item \textbf{Standardization (z-score):}
		$z=\dfrac{x-\mu}{\sigma}$ (zero mean, unit variance).
		\item \textbf{Min--max normalization:}
		$x'=\dfrac{x-x_{\min}}{x_{\max}-x_{\min}}\in[0,1]$.
	\end{itemize}
	Scaling is worked numerically in Chapter~\ref{ch:scaling}.
	
	%==============================================================================
	\chapter{The KNN Algorithm (Theory)}
	\label{ch:theory}
	%==============================================================================
	
	\section{The Algorithm}
	\begin{center}
		\fbox{\parbox{0.92\textwidth}{
				\textbf{$k$-Nearest Neighbours}
				\begin{enumerate}[nosep,leftmargin=2.2em]
					\item Choose the number of neighbours $k$ and a distance metric.
					\item For a query point $\vx_q$, compute the distance to \emph{every} training
					point.
					\item Sort the distances and select the $k$ nearest training points
					$N_k(\vx_q)$.
					\item \textbf{Classification:} predict the \emph{majority} class among
					$N_k$. \quad \textbf{Regression:} predict the \emph{average} target.
				\end{enumerate}
		}}
	\end{center}
	
	\section{The Decision Rule}
	\paragraph{Classification (majority vote).}
	\begin{equation}
		\hat y=\mode\{\,y_i : i\in N_k(\vx_q)\,\}
		=\argmax_{c}\ \sum_{i\in N_k(\vx_q)}\mathbf{1}(y_i=c),
		\label{eq:vote}
	\end{equation}
	where $\mathbf{1}(\cdot)$ is $1$ when true and $0$ otherwise.
	
	\paragraph{Regression (mean).}
	\begin{equation}
		\hat y=\frac{1}{k}\sum_{i\in N_k(\vx_q)} y_i .
		\label{eq:reg}
	\end{equation}
	
	\section{Distance-Weighted KNN}
	Closer neighbours are usually more relevant, so we can weight each vote by
	$w_i=1/d_i^{2}$ (or $1/d_i$):
	\begin{equation}
		\hat y_{\text{class}}=\argmax_{c}\!\!\sum_{i\in N_k,\,y_i=c}\!\! w_i,
		\qquad
		\hat y_{\text{reg}}=\frac{\sum_{i\in N_k} w_i\,y_i}{\sum_{i\in N_k} w_i}.
		\label{eq:weighted}
	\end{equation}
	Weighting also breaks ties naturally.
	
	\section{Choosing $k$ (the Bias--Variance Trade-off)}
	\begin{itemize}[nosep]
		\item \textbf{Small $k$} (e.g.\ $k=1$): very flexible, follows every point;
		\emph{low bias, high variance}---sensitive to noise and outliers.
		\item \textbf{Large $k$}: smoother boundary, averages over many points;
		\emph{high bias, low variance}---may blur real structure.
		\item A common rule of thumb is $k\approx\sqrt{m}$ (with $m$ training points),
		tuned by cross-validation.
		\item For binary classification pick an \textbf{odd} $k$ to avoid tied votes.
	\end{itemize}
	
	\section{Decision Boundary and Complexity}
	KNN induces a piecewise-linear decision boundary (for $k=1$, the boundaries of
	the \emph{Voronoi} cells of the training points). It has \textbf{no training
		cost} but each query costs $\mathcal{O}(mn)$ with brute force---slow for large
	$m$. Spatial indexes (Chapter~\ref{ch:types}) speed this up.
	
	%==============================================================================
	\chapter{Types and Variants of KNN}
	\label{ch:types}
	%==============================================================================
	
	\section{By Task and Weighting}
	\begin{description}[leftmargin=0pt,itemsep=3pt]
		\item[\textbf{KNN classification.}] Majority vote, equation \eqref{eq:vote}.
		\item[\textbf{KNN regression.}] Neighbour average, equation \eqref{eq:reg}.
		\item[\textbf{Distance-weighted KNN.}] Weight neighbours by $1/d^2$,
		equation \eqref{eq:weighted}.
		\item[\textbf{Radius-neighbours.}] Use \emph{all} points within a fixed radius
		$r$ instead of a fixed count $k$ (good for varying density).
	\end{description}
	
	\section{By Neighbour-Search Algorithm}
	\begin{center}
		\begin{tabular}{@{}lp{9.2cm}@{}}
			\toprule
			Method & Idea and cost\\
			\midrule
			Brute force & Compare the query to all $m$ points; $\mathcal{O}(mn)$ per query.
			Simple, exact, fine for small data.\\
			$k$-d tree & Binary space partition on feature axes; $\sim\mathcal{O}(\log m)$
			per query in low dimensions.\\
			Ball tree & Nested hyperspheres; better than $k$-d trees in higher dimensions.\\
			\bottomrule
		\end{tabular}
	\end{center}
	All three return the \emph{same} neighbours; they differ only in speed. In very
	high dimensions the ``curse of dimensionality'' makes all points nearly
	equidistant, and tree methods degrade toward brute force.
	
	%==============================================================================
	\chapter{Worked Example 1: Classification by Hand}
	\label{ch:classify}
	%==============================================================================
	
	\section{Problem Statement}
	Predict whether a student will \classpass{} or \classfail{} from two features:
	$x_1=$ hours studied, $x_2=$ attendance (\%).
	
	\section{Training Dataset}
	\begin{center}
		\renewcommand{\arraystretch}{1.2}
		\begin{tabular}{c c c c}
			\toprule
			\textbf{Student} & \textbf{Hours} & \textbf{Attendance (\%)} & \textbf{Class}\\
			$S_{j}$ & $h_{j}$ & $a_{j}$ & \\
			\midrule
			$S_{1}$  & 1 & 50 & \classfail\\
			$S_{2}$  & 2 & 55 & \classfail\\
			$S_{3}$  & 2 & 60 & \classfail\\
			$S_{4}$  & 3 & 60 & \classfail\\
			$S_{5}$  & 3 & 65 & \classpass\\
			$S_{6}$  & 4 & 65 & \classpass\\
			$S_{7}$  & 4 & 70 & \classpass\\
			$S_{8}$  & 5 & 70 & \classpass\\
			$S_{9}$  & 5 & 75 & \classpass\\
			$S_{10}$ & 6 & 80 & \classpass\\
			\bottomrule
		\end{tabular}
	\end{center}
	New student $X=(4,68)$; choose $k=3$. We use Euclidean distance
	\eqref{eq:euclid}.
	
	\section{Distance Calculation}
	For $X=(4,68)$ we apply $d(X,S_j)=\sqrt{(4-h_j)^2+(68-a_j)^2}$ to every student:
	\begin{align*}
		d(X,S_1)   &=\sqrt{(4-1)^2+(68-50)^2}=\sqrt{9+324}=\sqrt{333}\approx18.25,\\
		d(X,S_2)   &=\sqrt{(4-2)^2+(68-55)^2}=\sqrt{4+169}=\sqrt{173}\approx13.15,\\
		d(X,S_3)   &=\sqrt{(4-2)^2+(68-60)^2}=\sqrt{4+64}=\sqrt{68}\approx8.25,\\
		d(X,S_4)   &=\sqrt{(4-3)^2+(68-60)^2}=\sqrt{1+64}=\sqrt{65}\approx8.06,\\
		d(X,S_5)   &=\sqrt{(4-3)^2+(68-65)^2}=\sqrt{1+9}=\sqrt{10}\approx3.16,\\
		d(X,S_6)   &=\sqrt{(4-4)^2+(68-65)^2}=\sqrt{0+9}=\sqrt{9}=3.00,\\
		d(X,S_7)   &=\sqrt{(4-4)^2+(68-70)^2}=\sqrt{0+4}=\sqrt{4}=2.00,\\
		d(X,S_8)   &=\sqrt{(4-5)^2+(68-70)^2}=\sqrt{1+4}=\sqrt{5}\approx2.24,\\
		d(X,S_9)   &=\sqrt{(4-5)^2+(68-75)^2}=\sqrt{1+49}=\sqrt{50}\approx7.07,\\
		d(X,S_{10})&=\sqrt{(4-6)^2+(68-80)^2}=\sqrt{4+144}=\sqrt{148}\approx12.17.
	\end{align*}
	
	\begin{center}
		\renewcommand{\arraystretch}{1.2}
		\begin{tabular}{c c c c c}
			\toprule
			\textbf{Student} & \textbf{Hours} & \textbf{Attendance (\%)} & \textbf{Class} & $d(X,S_{j})$ \\
			$S_{j}$ & $h_{j}$ & $a_{j}$ & & $\sqrt{(4-h_j)^2+(68-a_j)^2}$ \\
			\midrule
			$S_{1}$  & 1 & 50 & \classfail & 18.25\\
			$S_{2}$  & 2 & 55 & \classfail & 13.15\\
			$S_{3}$  & 2 & 60 & \classfail & 8.25\\
			$S_{4}$  & 3 & 60 & \classfail & 8.06\\
			$S_{5}$  & 3 & 65 & \classpass & 3.16\\
			$S_{6}$  & 4 & 65 & \classpass & \textbf{3.00}\\
			$S_{7}$  & 4 & 70 & \classpass & \textbf{2.00}\\
			$S_{8}$  & 5 & 70 & \classpass & \textbf{2.24}\\
			$S_{9}$  & 5 & 75 & \classpass & 7.07\\
			$S_{10}$ & 6 & 80 & \classpass & 12.17\\
			\bottomrule
		\end{tabular}
	\end{center}
	The three smallest distances are $2.00,\ 2.24,\ 3.00$, so
	\[
	\boxed{\text{Nearest neighbours}=\{S_7,S_8,S_6\}.}
	\]
	
	\section{Majority Voting}
	\begin{center}
		\renewcommand{\arraystretch}{1.2}
		\begin{tabular}{c c c}
			\toprule
			\textbf{Neighbour} & \textbf{Distance} & \textbf{Class}\\
			\midrule
			S7 & 2.00 & \classpass\\
			S8 & 2.24 & \classpass\\
			S6 & 3.00 & \classpass\\
			\bottomrule
		\end{tabular}
	\end{center}
	Vote counts: $N(\text{Pass})=3$, $N(\text{Fail})=0$.
	\begin{resultbox}
		\[ \boxed{\hat y=\text{Pass}} \]
		The new student $X=(4,68)$ is classified as \textbf{Pass}.
	\end{resultbox}
	
	\section{Visual Interpretation}
	\begin{center}
		\begin{tikzpicture}
			\begin{axis}[
				width=0.82\textwidth,height=0.50\textwidth,
				xlabel={Hours studied}, ylabel={Attendance (\%)},
				xmin=0.5,xmax=6.5,ymin=47,ymax=82,grid=major,
				legend style={at={(0.02,0.98)},anchor=north west}]
				\addplot[only marks,mark=*,mark size=2.4pt,color=ccoral]
				coordinates {(1,50)(2,55)(2,60)(3,60)};
				\addlegendentry{Fail}
				\addplot[only marks,mark=*,mark size=2.4pt,color=cteal]
				coordinates {(3,65)(4,65)(4,70)(5,70)(5,75)(6,80)};
				\addlegendentry{Pass}
				\addplot[only marks,mark=*,mark size=4pt,color=corange]
				coordinates {(4,68)};
				\addlegendentry{New student}
				% circle the 3 nearest neighbours
				\addplot[only marks,mark=o,mark size=7pt,color=cnavy,thick]
				coordinates {(4,70)(5,70)(4,65)};
				\addlegendentry{$k=3$ neighbours}
			\end{axis}
		\end{tikzpicture}
	\end{center}
	
	\section{Effect of $k$}
	Because the two classes are cleanly separated here, the prediction is stable:
	\begin{center}
		\begin{tabular}{@{}cll@{}}
			\toprule
			$k$ & Neighbours & Prediction\\
			\midrule
			1 & $S_7$ & Pass\\
			3 & $S_7,S_8,S_6$ & Pass\\
			5 & $S_7,S_8,S_6,S_5,S_9$ & Pass\\
			7 & $S_7,S_8,S_6,S_5,S_9,S_4,S_3$ & Pass ($5$ vs $2$)\\
			\bottomrule
		\end{tabular}
	\end{center}
	In noisier data the choice of $k$ matters much more (Section~3.4): too small and
	a single mislabelled neighbour flips the answer; too large and distant,
	irrelevant points dilute the vote.
	
	%==============================================================================
	\chapter{Worked Example 2: KNN Regression by Hand}
	\label{ch:regress}
	%==============================================================================
	
	KNN also predicts \emph{numbers}. Here we estimate a house price (in lakh)
	from $x_1=$ number of rooms and $x_2=$ age (years). Query house $Q=(5,10)$,
	$k=3$.
	
	\section{Data and Distances}
	\begin{center}
		\renewcommand{\arraystretch}{1.2}
		\begin{tabular}{c c c c c}
			\toprule
			House & Rooms & Age & Price & $d(Q,\cdot)$\\
			\midrule
			H1 & 5 & 11 & 92 & $\sqrt{0+1}=\mathbf{1.00}$\\
			H2 & 6 & 11 & 88 & $\sqrt{1+1}=\mathbf{1.41}$\\
			H3 & 5 & 8  & 98 & $\sqrt{0+4}=\mathbf{2.00}$\\
			H4 & 7 & 12 & 80 & $\sqrt{4+4}=2.83$\\
			H5 & 2 & 9  & 70 & $\sqrt{9+1}=3.16$\\
			H6 & 8 & 14 & 74 & $\sqrt{9+16}=5.00$\\
			\bottomrule
		\end{tabular}
	\end{center}
	The $k=3$ nearest are $H_1,H_2,H_3$ (prices $92,88,98$).
	
	\section{Prediction}
	\paragraph{Simple average \eqref{eq:reg}:}
	\[
	\hat y=\frac{92+88+98}{3}=\frac{278}{3}\approx \boxed{92.67\ \text{lakh}.}
	\]
	\paragraph{Distance-weighted average \eqref{eq:weighted}} with $w_i=1/d_i^2$:
	\begin{center}
		\begin{tabular}{@{}ccccc@{}}
			\toprule
			House & $d$ & $d^2$ & $w=1/d^2$ & $w\cdot y$\\
			\midrule
			H1 & 1.00 & 1 & 1.000 & 92.00\\
			H2 & 1.41 & 2 & 0.500 & 44.00\\
			H3 & 2.00 & 4 & 0.250 & 24.50\\
			\midrule
			\multicolumn{3}{r}{$\sum$} & 1.750 & 160.50\\
			\bottomrule
		\end{tabular}
	\end{center}
	\[
	\hat y_{\text{weighted}}=\frac{160.50}{1.750}\approx \boxed{91.71\ \text{lakh}.}
	\]
	The weighted estimate leans toward the closest house $H_1$ (price $92$), as it
	should.
	
	\section{Including a Third Feature: Area (sq ft)}
	Suppose we now add $x_3=$ area (sq ft) to each house. The dataset and query
	become:
	\begin{center}
		\renewcommand{\arraystretch}{1.2}
		\begin{tabular}{c c c c c}
			\toprule
			House & Rooms & Age & Area (sq ft) & Price\\
			\midrule
			H1 & 5 & 11 & 1450 & 92\\
			H2 & 6 & 11 & 1600 & 88\\
			H3 & 5 & 8  & 1200 & 98\\
			H4 & 7 & 12 & 1520 & 80\\
			H5 & 2 & 9  & 900  & 70\\
			H6 & 8 & 14 & 2000 & 74\\
			\bottomrule
		\end{tabular}
	\end{center}
	Query $Q=(5,\,10,\,1500)$. The Euclidean distance now runs over three features:
	\[
	d(Q,H)=\sqrt{(\Delta\text{rooms})^2+(\Delta\text{age})^2+(\Delta\text{area})^2}.
	\]
	
	\subsection*{Unscaled distances}
	\begin{center}
		\renewcommand{\arraystretch}{1.2}
		\begin{tabular}{c c c c}
			\toprule
			House & $(\Delta r)^2+(\Delta a)^2+(\Delta\text{area})^2$ & $d^2$ & $d$\\
			\midrule
			H1 & $0+1+2500$    & $2501$   & $\mathbf{50.01}$\\
			H2 & $1+1+10000$   & $10002$  & $\mathbf{100.01}$\\
			H3 & $0+4+90000$   & $90004$  & $300.01$\\
			H4 & $4+4+400$     & $408$    & $\mathbf{20.20}$\\
			H5 & $9+1+360000$  & $360010$ & $600.01$\\
			H6 & $9+16+250000$ & $250025$ & $500.02$\\
			\bottomrule
		\end{tabular}
	\end{center}
	The three nearest are now $\{H_4,H_1,H_2\}$ (prices $80,92,88$)---\emph{not}
	$\{H_1,H_2,H_3\}$ as before. The simple prediction is
	\[
	\hat y=\frac{80+92+88}{3}=\frac{260}{3}\approx\boxed{86.67\ \text{lakh},}
	\]
	and the distance-weighted prediction ($w=1/d^2$) is
	\begin{center}
		\begin{tabular}{@{}ccccc@{}}
			\toprule
			House & $d^2$ & $w=1/d^2$ & Price & $w\cdot y$\\
			\midrule
			H4 & $408$   & $0.002451$ & $80$ & $0.1961$\\
			H1 & $2501$  & $0.000400$ & $92$ & $0.0368$\\
			H2 & $10002$ & $0.000100$ & $88$ & $0.0088$\\
			\midrule
			\multicolumn{2}{r}{$\sum$} & $0.002951$ & & $0.2417$\\
			\bottomrule
		\end{tabular}
	\end{center}
	\[
	\hat y_{\text{weighted}}=\frac{0.2417}{0.002951}\approx\boxed{81.90\ \text{lakh}.}
	\]
	
	\begin{conceptbox}
		\textbf{Why did the answer change?} Area ranges in the \emph{thousands} while
		rooms and age are single/double digits, so the term $(\Delta\text{area})^2$
		dominates $d^2$ completely---distances are essentially the area gap alone. The
		prediction dropped from $92.67$ to $86.67$ \emph{not} because area is truly more
		informative, but purely because of its larger units. This is exactly the
		feature-scaling problem of Chapter~\ref{ch:scaling}.
	\end{conceptbox}
	
	\subsection*{Standardized distances (the fair comparison)}
	Standardize each feature over the six houses, $z=(x-\mu)/\sigma$, with
	\[
	\mu=(5.5,\,10.833,\,1445),\qquad \sigma=(1.893,\,1.951,\,340.28).
	\]
	The query becomes $Q_z=(-0.264,-0.427,\,0.162)$, and the scaled distances are:
	\begin{center}
		\renewcommand{\arraystretch}{1.2}
		\begin{tabular}{c c c}
			\toprule
			House & Price & Scaled $d$\\
			\midrule
			H1 & 92 & $\mathbf{0.533}$\\
			H2 & 88 & $\mathbf{0.793}$\\
			H3 & 98 & $\mathbf{1.352}$\\
			H4 & 80 & $1.473$\\
			H5 & 70 & $2.426$\\
			H6 & 74 & $2.979$\\
			\bottomrule
		\end{tabular}
	\end{center}
	With every feature on an equal footing the three nearest are again
	$\{H_1,H_2,H_3\}$, giving
	\[
	\hat y=\frac{92+88+98}{3}\approx 92.67\ \text{lakh}\quad(\text{weighted}\approx 91.45\ \text{lakh}),
	\]
	the same as the two-feature answer---area now contributes without dominating.
	
	\subsection*{Summary}
	\begin{center}
		\renewcommand{\arraystretch}{1.2}
		\begin{tabular}{l c c c}
			\toprule
			Scenario & Neighbours & Simple avg & Weighted avg\\
			\midrule
			2 features (rooms, age)      & $H_1,H_2,H_3$ & $92.67$ & $91.71$\\
			3 features, \textbf{unscaled}  & $H_4,H_1,H_2$ & $86.67$ & $81.90$\\
			3 features, \textbf{standardized} & $H_1,H_2,H_3$ & $92.67$ & $91.45$\\
			\bottomrule
		\end{tabular}
	\end{center}
	Adding a feature can change the prediction dramatically---but if the new feature
	has a larger scale, that change is an \emph{artefact of units}. Always scale
	before adding features of different ranges.
	
	%==============================================================================
	\chapter{Worked Example 3: Feature Scaling by Hand}
	\label{ch:scaling}
	%==============================================================================
	
	We revisit the classification data of Chapter~\ref{ch:classify} to see
	\emph{why} scaling (Section~\ref{sec:scaling}) matters. In the raw distances,
	attendance (range $50$--$80$) completely dominates hours (range $1$--$6$): every
	distance was essentially the attendance gap. We now standardize both features.
	
	\section{Training Dataset}
	\begin{center}
		\renewcommand{\arraystretch}{1.2}
		\begin{tabular}{c c c c}
			\toprule
			\textbf{Student} & \textbf{Hours} & \textbf{Attendance (\%)} & \textbf{Class}\\
			$S_{j}$ & $h_{j}$ & $a_{j}$ & \\
			\midrule
			$S_{1}$  & 1 & 50 & \classfail\\
			$S_{2}$  & 2 & 55 & \classfail\\
			$S_{3}$  & 2 & 60 & \classfail\\
			$S_{4}$  & 3 & 60 & \classfail\\
			$S_{5}$  & 3 & 65 & \classpass\\
			$S_{6}$  & 4 & 65 & \classpass\\
			$S_{7}$  & 4 & 70 & \classpass\\
			$S_{8}$  & 5 & 70 & \classpass\\
			$S_{9}$  & 5 & 75 & \classpass\\
			$S_{10}$ & 6 & 80 & \classpass\\
			\bottomrule
		\end{tabular}
	\end{center}
	
	\section{Compute the Feature Statistics}
	Over the ten students,
	\[
	\bar{h}=\sum\limits_{j=1}^{10} h_{j}=3.5,\quad \sigma_{\text{h}}=\sqrt{\dfrac{1}{10}\sum\limits_{j=1}^{10}\left( h_j-\bar{h} \right)^2}=1.5
	\]
	
	\[
	\bar{a}=\sum\limits_{j=1}^{10} a_{j}=65,\quad \sigma_{\text{a}}=\sqrt{\dfrac{1}{10}\sum\limits_{j=1}^{10}\left( a_j-\bar{a} \right)^2}\approx8.66 .
	\]
	Standardize with $z=(x-\mu)/\sigma$. For the query $X=(4,68)$:
	\[
	z_1=\frac{4-3.5}{1.5}=0.333,\qquad z_2=\frac{68-65}{8.66}=0.346 .
	\]
	
	\section{Standardized Distances}
	For example, $S_6=(4,65)\to(0.333,0)$ and $S_8=(5,70)\to(1.0,0.577)$, giving
	\[
	d(X,S_6)=\sqrt{0^2+0.346^2}=0.346,\qquad
	d(X,S_8)=\sqrt{0.667^2+0.231^2}=0.706 .
	\]
	\begin{center}
		\small
		\setlength{\tabcolsep}{5pt}
		\renewcommand{\arraystretch}{1.25}
		\begin{tabular}{c| cc | cc | c | cc}
			\toprule
			\multirow{2}{*}{\shortstack{Student\\ $S_j$}}
			& \multicolumn{2}{c|}{Hours, $h_j$}
			& \multicolumn{2}{c|}{Attendance (\%)}
			& \multirow{2}{*}{Class}
			& \multicolumn{2}{c}{$d(X,S_j)$}\\
			\cmidrule(lr){2-3}\cmidrule(lr){4-5}\cmidrule(lr){7-8}
			& Raw & \shortstack{Standardized\\[2pt]$\dfrac{h_j-\bar{h}}{\sigma_h}$}
			& Raw & \shortstack{Standardized\\[2pt]$\dfrac{a_j-\bar{a}}{\sigma_a}$}
			& & Raw & Scaled\\
			\midrule
			$S_1$  & 1 & $-1.667$ & 50 & $-1.732$ & \classfail & 18.25 & 2.88\\
			$S_2$  & 2 & $-1.000$ & 55 & $-1.155$ & \classfail & 13.15 & 2.01\\
			$S_3$  & 2 & $-1.000$ & 60 & $-0.577$ & \classfail & 8.25  & 1.62\\
			$S_4$  & 3 & $-0.333$ & 60 & $-0.577$ & \classfail & 8.06  & 1.14\\
			$S_5$  & 3 & $-0.333$ & 65 & $0.000$  & \classpass & 3.16  & 0.75\\
			$S_6$  & 4 & $0.333$  & 65 & $0.000$  & \classpass & 3.00  & 0.35\\
			$S_7$  & 4 & $0.333$  & 70 & $0.577$  & \classpass & 2.00  & 0.23\\
			$S_8$  & 5 & $1.000$  & 70 & $0.577$  & \classpass & 2.24  & 0.71\\
			$S_9$  & 5 & $1.000$  & 75 & $1.155$  & \classpass & 7.07  & 1.05\\
			$S_{10}$ & 6 & $1.667$  & 80 & $1.732$  & \classpass & 12.17 & 1.92\\
			\bottomrule
		\end{tabular}
	\end{center}
	
	\begin{resultbox}
		Under \emph{raw} distances the $3$ nearest were $\{S_7,S_8,S_6\}$; after
		\emph{scaling} they are $\{S_7,S_6,S_8\}$---$S_6$ and $S_8$ have \textbf{swapped
			order} because standardization restored the influence of ``hours''. Here the
		class (Pass) is unchanged, but on many datasets scaling changes the neighbour
		set and therefore the prediction. Always scale in practice.
	\end{resultbox}
	
	%==============================================================================
	\chapter{Worked Example 4: Min--Max Normalization by Hand}
	\label{ch:minmax}
	%==============================================================================
	
	Chapter~\ref{ch:scaling} used \emph{standardization} (the $z$-score) to put the
	two features on a common scale. \emph{Min--max normalization} is the other
	standard rescaling: instead of centring on the mean, it stretches each feature
	so that its smallest value becomes $0$ and its largest becomes $1$.
	
	\section{The Min--Max Formula}
	For a feature $x$ with minimum $x_{\min}$ and maximum $x_{\max}$,
	\begin{equation}
	x'=\frac{x-x_{\min}}{x_{\max}-x_{\min}}\ \in[0,1],
	\label{eq:minmax}
	\end{equation}
	so $x'=0$ at the feature's minimum, $x'=1$ at its maximum, and everything else
	lands proportionally in between.
	
	\section{Feature Ranges}
	Using the same ten students of Chapter~\ref{ch:scaling}:
	\[
	h_{\min}=1,\ h_{\max}=6\ \Rightarrow\ h'=\frac{h-1}{5};\qquad
	a_{\min}=50,\ a_{\max}=80\ \Rightarrow\ a'=\frac{a-50}{30}.
	\]
	The query $X=(4,68)$ therefore normalizes to
	\[
	h'=\frac{4-1}{5}=0.6,\qquad a'=\frac{68-50}{30}=\frac{18}{30}=0.6,
	\qquad\text{so } X'=(0.6,\,0.6).
	\]
	
	\section{Normalized Distances}
	For example, $S_6=(4,65)\to(0.600,0.500)$ and $S_7=(4,70)\to(0.600,0.667)$,
	giving
	\[
	d(X,S_6)=\sqrt{0^2+0.100^2}=0.100,\qquad
	d(X,S_7)=\sqrt{0^2+0.067^2}=0.067 .
	\]
	\begin{center}
	\small
	\setlength{\tabcolsep}{5pt}
	\renewcommand{\arraystretch}{1.25}
	\begin{tabular}{c| cc | cc | c | cc}
	\toprule
	\multirow{2}{*}{\shortstack{Student\\ $S_j$}}
	& \multicolumn{2}{c|}{Hours, $h_j$}
	& \multicolumn{2}{c|}{Attendance (\%)}
	& \multirow{2}{*}{Class}
	& \multicolumn{2}{c}{$d(X,S_j)$}\\
	\cmidrule(lr){2-3}\cmidrule(lr){4-5}\cmidrule(lr){7-8}
	& Raw & \shortstack{Normalized\\[2pt]$\dfrac{h_j-h_{\min}}{h_{\max}-h_{\min}}$}
	& Raw & \shortstack{Normalized\\[2pt]$\dfrac{a_j-a_{\min}}{a_{\max}-a_{\min}}$}
	& & Raw & Norm.\\
	\midrule
	$S_1$  & 1 & $0.000$ & 50 & $0.000$ & \classfail & 18.25 & 0.85\\
	$S_2$  & 2 & $0.200$ & 55 & $0.167$ & \classfail & 13.15 & 0.59\\
	$S_3$  & 2 & $0.200$ & 60 & $0.333$ & \classfail & 8.25  & 0.48\\
	$S_4$  & 3 & $0.400$ & 60 & $0.333$ & \classfail & 8.06  & 0.33\\
	$S_5$  & 3 & $0.400$ & 65 & $0.500$ & \classpass & 3.16  & 0.22\\
	$S_6$  & 4 & $0.600$ & 65 & $0.500$ & \classpass & 3.00  & 0.10\\
	$S_7$  & 4 & $0.600$ & 70 & $0.667$ & \classpass & 2.00  & 0.07\\
	$S_8$  & 5 & $0.800$ & 70 & $0.667$ & \classpass & 2.24  & 0.21\\
	$S_9$  & 5 & $0.800$ & 75 & $0.833$ & \classpass & 7.07  & 0.31\\
	$S_{10}$ & 6 & $1.000$ & 80 & $1.000$ & \classpass & 12.17 & 0.57\\
	\bottomrule
	\end{tabular}
	\end{center}
	
	\begin{resultbox}
	The three smallest normalized distances are $0.07\ (S_7),\ 0.10\ (S_6),\
	0.21\ (S_8)$, so the $k=3$ neighbours are $\{S_7,S_6,S_8\}$---all
	\classpass{}---and $\hat y=\text{Pass}$. As with standardization, $S_6$ now
	ranks ahead of $S_8$ (they were $S_8$ before $S_6$ under the raw distances),
	because rescaling gives ``hours'' a fair say.
	\end{resultbox}
	
	\section{Standardization vs.\ Min--Max Normalization}
	Both rescalings fix the ``attendance dominates'' problem; they differ in how.
	\begin{center}
	\begin{tabular}{@{}p{6.4cm}p{6.4cm}@{}}
	\toprule
	\textbf{Standardization ($z$-score)} & \textbf{Min--max normalization}\\
	\midrule
	$x'=\dfrac{x-\mu}{\sigma}$; mean $0$, variance $1$. &
	$x'=\dfrac{x-x_{\min}}{x_{\max}-x_{\min}}$; range $[0,1]$.\\
	Unbounded; centred on the mean. & Bounded to $[0,1]$; anchored to the extremes.\\
	More robust to outliers (uses $\mu,\sigma$). & Sensitive to outliers (a single
	extreme value sets $x_{\min}$ or $x_{\max}$).\\
	Good default for distance-based methods. & Handy when a bounded range is needed
	(e.g.\ image pixels, neural-net inputs).\\
	\bottomrule
	\end{tabular}
	\end{center}
	On this dataset both give the identical neighbour set $\{S_7,S_6,S_8\}$ and the
	prediction \textbf{Pass}; on data with outliers they can disagree.
	
	%==============================================================================
	\chapter{Worked Example 5: Multi-Class Classification (Three Classes)}
	\label{ch:multiclass}
	%==============================================================================
	
	KNN needs \emph{no} change to handle more than two classes: it still takes the
	$k$ nearest neighbours and predicts the majority class. Here a query point $P$
	sits amid three classes A, B and C (Figure~\ref{fig:multiclass}).
	
	\section{The Data}
	Five labelled points per class:
	\begin{center}
	\renewcommand{\arraystretch}{1.2}
	\begin{tabular}{c c c}
	\toprule
	\textbf{Class A} & \textbf{Class B} & \textbf{Class C}\\
	\midrule
	$(1,6)$ & $(8,7)$ & $(4,1)$\\
	$(2,7)$ & $(9,6)$ & $(5,2)$\\
	$(2,5)$ & $(9,8)$ & $(6,1)$\\
	$(3,7)$ & $(8,5)$ & $(3,2)$\\
	$(1,8)$ & $(7,7)$ & $(6,3)$\\
	\bottomrule
	\end{tabular}
	\end{center}
	Query point $P=(6,5)$; use $k=7$ and Euclidean distance.
	
	\section{Distances from $P$}
	Squared distance is enough to rank the neighbours ($d^2=(6-x)^2+(5-y)^2$). The
	fifteen points are listed \textbf{by class} (A, then B, then C); the
	\textbf{seven nearest} (Rank $\le 7$) are shown in bold.
	\begin{center}
	\small
	\renewcommand{\arraystretch}{1.15}
	\begin{tabular}{c c c c c}
	\toprule
	Class & Point & $d^2$ & $d$ & Rank\\
	\midrule
	\multirow{5}{*}{\textbf{A}}
	 & $(1,6)$ & $26$ & $5.10$          & $14$\\
	 & $(2,7)$ & $20$ & $4.47$          & $12$\\
	 & $(2,5)$ & $16$ & $4.00$          & $8$\\
	 & $(3,7)$ & $13$ & $\mathbf{3.61}$ & $\mathbf{7}$\\
	 & $(1,8)$ & $34$ & $5.83$          & $15$\\
	\midrule
	\multirow{5}{*}{\textbf{B}}
	 & $(8,7)$ & $8$  & $\mathbf{2.83}$ & $\mathbf{4}$\\
	 & $(9,6)$ & $10$ & $\mathbf{3.16}$ & $\mathbf{6}$\\
	 & $(9,8)$ & $18$ & $4.24$          & $11$\\
	 & $(8,5)$ & $4$  & $\mathbf{2.00}$ & $\mathbf{2}$\\
	 & $(7,7)$ & $5$  & $\mathbf{2.24}$ & $\mathbf{3}$\\
	\midrule
	\multirow{5}{*}{\textbf{C}}
	 & $(4,1)$ & $20$ & $4.47$          & $13$\\
	 & $(5,2)$ & $10$ & $\mathbf{3.16}$ & $\mathbf{5}$\\
	 & $(6,1)$ & $16$ & $4.00$          & $9$\\
	 & $(3,2)$ & $18$ & $4.24$          & $10$\\
	 & $(6,3)$ & $4$  & $\mathbf{2.00}$ & $\mathbf{1}$\\
	\bottomrule
	\end{tabular}
	\end{center}
	
	\section{Majority Vote}
	Counting classes among the seven nearest:
	\[
	N(\text{B})=4\ \{(8,5),(7,7),(8,7),(9,6)\},\quad
	N(\text{C})=2\ \{(6,3),(5,2)\},\quad
	N(\text{A})=1\ \{(3,7)\}.
	\]
	\begin{resultbox}
	\[ \boxed{\hat y=\text{Class B}} \]
	Class B has the most votes ($4$ of $7$), so $P$ is assigned to \textbf{Class B}.
	\end{resultbox}
	\begin{remark}
	The single closest point is a \emph{tie} between a C point $(6,3)$ and a B point
	$(8,5)$ (both $d=2$), so $1$-NN is ambiguous---yet the $k=7$ vote is decisively
	B. In fact B wins at $k=3$ ($2$--$1$), $k=5$ ($3$--$2$) and $k=7$ ($4$--$2$--$1$)
	alike. Majority voting is more robust than trusting the single nearest point.
	\end{remark}
	
	\section{Visual Interpretation}
	\begin{figure}[ht]
	\centering
	\begin{tikzpicture}
	\begin{axis}[
	  width=0.82\textwidth, unit vector ratio=1 1,
	  axis lines=middle, xlabel={$x_1$}, ylabel={$x_2$},
	  xlabel style={anchor=north west}, ylabel style={anchor=south east},
	  xmin=0, xmax=10.5, ymin=0, ymax=9.6,
	  xtick={1,2,3,4,5,6,7,8,9,10}, ytick={1,2,3,4,5,6,7,8,9},
	  tick label style={font=\scriptsize}, grid=both,
	  major grid style={line width=0.2pt,draw=gray!30},
	  legend pos=north west, legend cell align=left, clip=false]
	% dashed class clouds
	\draw[cnavy,dashed,thick,fill=cnavy!8]   (axis cs:1.8,6.6) ellipse [x radius=1.6,y radius=1.9];
	\draw[cteal,dashed,thick,fill=cteal!10]  (axis cs:8.2,6.6) ellipse [x radius=1.5,y radius=1.9];
	\draw[ccoral,dashed,thick,fill=ccoral!8] (axis cs:4.5,1.8) ellipse [x radius=2.1,y radius=1.4];
	% data points (+ legend)
	\addplot[only marks,mark=*,mark size=2.5pt,color=cnavy] coordinates {(1,6)(2,7)(2,5)(3,7)(1,8)};
	\addlegendentry{Class A}
	\addplot[only marks,mark=square*,mark size=2.5pt,color=cteal] coordinates {(8,7)(9,6)(9,8)(8,5)(7,7)};
	\addlegendentry{Class B}
	\addplot[only marks,mark=diamond*,mark size=3pt,color=ccoral] coordinates {(4,1)(5,2)(6,1)(3,2)(6,3)};
	\addlegendentry{Class C}
	\addplot[only marks,mark=star,mark size=4pt,color=corange] coordinates {(6,5)};
	\addlegendentry{Query $P$}
	% arrows from P to the nearest member of each class
	\draw[-{Latex[length=2mm]},thick] (axis cs:6,5)--(axis cs:3,7);
	\draw[-{Latex[length=2mm]},thick] (axis cs:6,5)--(axis cs:8,5);
	\draw[-{Latex[length=2mm]},thick] (axis cs:6,5)--(axis cs:6,3);
	% coordinate labels
	\node[font=\scriptsize,anchor=south] at (axis cs:1,6) {$(1,6)$};
	\node[font=\scriptsize,anchor=south] at (axis cs:2,7) {$(2,7)$};
	\node[font=\scriptsize,anchor=south] at (axis cs:2,5) {$(2,5)$};
	\node[font=\scriptsize,anchor=south] at (axis cs:3,7) {$(3,7)$};
	\node[font=\scriptsize,anchor=south] at (axis cs:1,8) {$(1,8)$};
	\node[font=\scriptsize,anchor=south] at (axis cs:8,7) {$(8,7)$};
	\node[font=\scriptsize,anchor=south] at (axis cs:9,6) {$(9,6)$};
	\node[font=\scriptsize,anchor=south] at (axis cs:9,8) {$(9,8)$};
	\node[font=\scriptsize,anchor=south] at (axis cs:8,5) {$(8,5)$};
	\node[font=\scriptsize,anchor=south] at (axis cs:7,7) {$(7,7)$};
	\node[font=\scriptsize,anchor=south] at (axis cs:4,1) {$(4,1)$};
	\node[font=\scriptsize,anchor=south] at (axis cs:5,2) {$(5,2)$};
	\node[font=\scriptsize,anchor=south] at (axis cs:6,1) {$(6,1)$};
	\node[font=\scriptsize,anchor=south] at (axis cs:3,2) {$(3,2)$};
	\node[font=\scriptsize,anchor=west] at (axis cs:6,3) {$(6,3)$};
	\node[font=\scriptsize\bfseries,anchor=north west] at (axis cs:6,5) {$P(6,5)$};
	\end{axis}
	\end{tikzpicture}
	\caption{Three-class KNN. The dashed clouds show classes A, B and C; arrows
	point from the query $P$ to the nearest member of each class. Although Class C
	owns one of the two joint-nearest points, the neighbourhood of $P$ holds more
	Class B points, so the $k=7$ majority vote returns \textbf{Class B}.}
	\label{fig:multiclass}
	\end{figure}
	
	%==============================================================================
	\chapter{Advantages, Disadvantages and Code}
	%==============================================================================
	
	\section{Pros and Cons}
	\begin{center}
		\begin{tabular}{@{}p{6.4cm}p{6.4cm}@{}}
			\toprule
			\textbf{Advantages} & \textbf{Disadvantages}\\
			\midrule
			Simple, intuitive, no training phase. & Slow at prediction: $\mathcal{O}(mn)$
			per query (must scan all data).\\
			Naturally handles multi-class problems. & Stores the entire dataset (high
			memory).\\
			No assumptions about data distribution (non-parametric). & Sensitive to feature
			scaling and to irrelevant features.\\
			Adapts instantly to new data (just add points). & Degrades in high dimensions
			(curse of dimensionality); must choose $k$.\\
			\bottomrule
		\end{tabular}
	\end{center}
	
	\section{Board Summary}
	\[
	\boxed{\;k=3\ \rightarrow\ \{S_7,S_8,S_6\}\ \rightarrow\
		\{\text{Pass},\text{Pass},\text{Pass}\}\ \rightarrow\ \text{Pass}\;}
	\]
	
	\section{KNN in Python (scikit-learn)}
	\begin{lstlisting}[style=python,caption={Classification and regression with scikit-learn.}]
		import numpy as np
		from sklearn.neighbors import KNeighborsClassifier
		from sklearn.preprocessing import StandardScaler
		
		X = np.array([[1,50],[2,55],[2,60],[3,60],[3,65],
		[4,65],[4,70],[5,70],[5,75],[6,80]], dtype=float)
		y = np.array([0,0,0,0,1,1,1,1,1,1])   # 0=Fail, 1=Pass
		
		# scale features, then fit KNN with k=3
		Xs = StandardScaler().fit_transform(X)
		knn = KNeighborsClassifier(n_neighbors=3).fit(Xs, y)
		
		q = StandardScaler().fit(X).transform([[4,68]])
		print("prediction:", knn.predict(q))     # -> [1] = Pass
		
		# distance-weighted variant:
		# KNeighborsClassifier(n_neighbors=3, weights='distance')
	\end{lstlisting}
	
	%==============================================================================
	\chapter{Exercises}
	%==============================================================================
	
	\section{Solved Problems}
	\begin{example}[One distance]
		Compute $d(X,S_3)$ for $X=(4,68)$, $S_3=(2,60)$.\\
		\textbf{Solution.} $\sqrt{(4-2)^2+(68-60)^2}=\sqrt{4+64}=\sqrt{68}\approx8.25.$
	\end{example}
	\begin{example}[$k=1$]
		Using the data of Chapter~\ref{ch:classify}, classify $X=(4,68)$ with $k=1$.\\
		\textbf{Solution.} The single nearest point is $S_7$ ($d=2.00$, Pass), so
		$\hat y=$ Pass.
	\end{example}
	\begin{example}[Regression average]
		Neighbours have targets $10,14,18$. Predict with $k=3$ (unweighted).\\
		\textbf{Solution.} $\hat y=(10+14+18)/3=14.$
	\end{example}
	\begin{example}[Why odd $k$?]
		Why choose an odd $k$ for two-class problems?\\
		\textbf{Solution.} An odd number of votes cannot split evenly between two
		classes, so majority voting never ties.
	\end{example}
	
	\section{Unsolved Problems}
	\begin{exercise}[Manhattan distance]
		Recompute the three nearest neighbours of $X=(4,68)$ using Manhattan ($L_1$)
		distance instead of Euclidean. Does the prediction change?
	\end{exercise}
	\begin{exercise}[Weighted vote]
		For the neighbours $S_7(2.00),S_8(2.24),S_6(3.00)$ (all Pass) and a hypothetical
		far Fail neighbour at $d=1.5$, redo the vote with weights $w=1/d^2$. Which class
		wins?
	\end{exercise}
	\begin{exercise}[Choosing $k$]
		With $m=10$ training points, what does the rule of thumb $k\approx\sqrt{m}$
		suggest? Discuss the trade-off if you instead used $k=1$ or $k=9$.
	\end{exercise}
	\begin{exercise}[Scaling]
		Standardize the query $X=(4,68)$ and point $S_5=(3,65)$ and verify the scaled
		distance $d\approx0.75$.
	\end{exercise}
	\begin{exercise}[Regression, weighted]
		For the houses of Chapter~\ref{ch:regress}, predict the price with $k=2$
		(unweighted) and compare with the $k=3$ answer.
	\end{exercise}
	\begin{exercise}[Curse of dimensionality]
		Explain, in one paragraph, why KNN struggles when the number of features $n$ is
		very large even if $m$ is fixed.
	\end{exercise}
	
	%==============================================================================
	\appendix
	\chapter{Notation Summary}
	%==============================================================================
	\begin{center}
		\begin{tabular}{@{}ll@{}}
			\toprule
			Symbol & Meaning\\
			\midrule
			$m$ & number of training points\\
			$n$ & number of features\\
			$k$ & number of neighbours used\\
			$\vx_q$ & query (new) point to be classified\\
			$N_k(\vx_q)$ & set of the $k$ nearest training points\\
			$d(\vx,\vy)$ & distance between two points\\
			$w_i=1/d_i^2$ & distance weight of neighbour $i$\\
			$\hat y$ & predicted label (class) or value (regression)\\
			\bottomrule
		\end{tabular}
	\end{center}
	
\end{document}
